\section{Block Schur Complements}\label{sec:schur_complement}

\begin{theorem}[Wolf's printed block Schur complement statement]\label{thm:schur_tfae}
  \notready
  \uses{def:schurComplement, def:blockMatrix}
  Let $P \in \MN{D_1}$, $R \in \MN{D_2}$ be positive semi-definite
  and let $Q \in \MN{D_1,D_2}$. Wolf's Theorem~5.2 states that the
  following are equivalent:
  \begin{enumerate}
    \item The block matrix $\begin{pmatrix} P & Q \\ Q^\dagger & R \end{pmatrix}$ is PSD,
    \item $\ker(R) \subseteq \ker(Q)$ and
      $P \ge Q R^+ Q^\dagger$,
    \item $\ker(R) \subseteq \ker(Q)$ and
      $\|P^{-1/2}Q R^{-1/2}\|_\infty \le 1$.
  \end{enumerate}
  Here $R^+$ denotes the pseudoinverse (inverse on the support) and the
  inverse square roots are extended by zero on the respective kernels.

  \textbf{Local fix (Wolf Theorem 5.2, condition (3)):}
  The third condition is false as printed when $P$ is singular. It becomes
  equivalent to the first two after adjoining
  $\ker(P)\subseteq\ker(Q^\dagger)$. The counterexample and corrected theorem
  are recorded below and in
  \texttt{docs/paper-gaps/schur\_complement\_tfae.tex}.
\end{theorem}

\begin{definition}[Block matrix]\label{def:blockMatrix}
  \lean{SchurComplement.blockMatrix}
  \leanok
  A $2 \times 2$ block matrix $\begin{pmatrix} P & Q \\ Q^\dagger & R \end{pmatrix}$
  indexed by $\Fin{D_1} \oplus \Fin{D_2}$.
\end{definition}

\begin{lemma}[Hermiticity of the block matrix]\label{lem:blockMatrix_isHermitian}
  \lean{SchurComplement.blockMatrix_isHermitian}
  \leanok
  \uses{def:blockMatrix}
  If $P$ and $R$ are Hermitian, then the block matrix
  $\begin{pmatrix} P & Q \\ Q^\dagger & R \end{pmatrix}$ is Hermitian.
\end{lemma}

\begin{theorem}[Quadratic form of the block matrix]
  \label{thm:block_quadratic_form}
  \lean{SchurComplement.block_quadratic_form}
  \leanok
  \uses{def:blockMatrix}
  For vectors $x\in\mathbb C^{D_1}$ and $y\in\mathbb C^{D_2}$,
  \begin{align}
    \begin{pmatrix}x\\y\end{pmatrix}^{\!\dagger}
      \begin{pmatrix}P&Q\\Q^\dagger&R\end{pmatrix}
      \begin{pmatrix}x\\y\end{pmatrix}
    =x^\dagger P x+x^\dagger Qy+y^\dagger Q^\dagger x+y^\dagger Ry.
  \end{align}
\end{theorem}

\begin{proof}\leanok
  Multiplication by the block matrix sends $(x,y)$ to
  $(Px+Qy,Q^\dagger x+Ry)$. Taking the inner product with $(x,y)$ and
  distributing over both sums gives the four displayed terms.
\end{proof}

\begin{definition}[Schur complement]\label{def:schurComplement}
  \lean{SchurComplement.schurComplement}
  \leanok
  \uses{def:pinv}
  The pseudoinverse Schur complement $P - Q R^+ Q^\dagger$.
\end{definition}

\begin{theorem}[Support absorption under kernel inclusion]
  \label{thm:support_absorption_of_kernel_inclusion}
  \lean{Matrix.IsHermitian.mul_supportProj_eq_self_of_mulVec_kernel_le}
  \leanok
  \uses{def:hermitian_support_projection}
  Let $R$ be Hermitian and let $Q$ be a matrix with the same column index
  such that $\ker(R)\subseteq\ker(Q)$. Then
  $Q\operatorname{supp}(R)=Q$.
\end{theorem}

\begin{proof}\leanok
  Since $R(1-\operatorname{supp}(R))=0$, the kernel inclusion gives
  $Q(1-\operatorname{supp}(R))=0$. Expanding this equality yields
  $Q-Q\operatorname{supp}(R)=0$, and hence
  $Q\operatorname{supp}(R)=Q$.
\end{proof}

\begin{theorem}[Pseudoinverse support identities]
  \label{thm:schur_pinv_support_identities}
  \lean{SchurComplement.R_mul_pinv_eq_supportProj}
  \lean{SchurComplement.pinv_mul_self_eq_supportProj}
  \lean{SchurComplement.supportProj_mul_pinv_eq_pinv}
  \lean{SchurComplement.pinv_isHermitian}
  \leanok
  \uses{def:pinv, def:hermitian_support_projection}
  If $R$ is positive semi-definite, then
  $RR^+=R^+R=\operatorname{supp}(R)$,
  $\operatorname{supp}(R)R^+=R^+$, and $R^+$ is Hermitian.
\end{theorem}

\begin{proof}\leanok
  \uses{lem:mul_pinv_eq_supportProj}
  The pseudoinverse-factorization identity, together with equality of the
  supports of $R^2$ and $R$, gives $RR^+=\operatorname{supp}(R)$.
  The spectral calculus shows that the support inverse of $R^2$ commutes with
  $R$; the definition of $R^+$ then shows that $R^+$ is Hermitian. Taking
  adjoints in the first identity gives
  $R^+R=\operatorname{supp}(R)$. Finally, write $R=\sqrt R\sqrt R$.
  The support projection absorbs $\sqrt R$, and the defining factorization of
  $R^+$ therefore gives $\operatorname{supp}(R)R^+=R^+$.
\end{proof}

\begin{theorem}[Kernel inclusion from block positivity]
  \label{thm:block_psd_kernel_inclusion}
  \lean{SchurComplement.ker_subset_of_block_psd}
  \leanok
  \uses{def:blockMatrix}
  If $\begin{pmatrix}P&Q\\Q^\dagger&R\end{pmatrix}$ is positive
  semi-definite, then $\ker(R)\subseteq\ker(Q)$.
\end{theorem}

\begin{proof}\leanok
  \uses{thm:block_quadratic_form}
  A vector in $\ker(R)$ has zero quadratic form in the lower-right principal
  subspace. Positivity therefore places the corresponding block vector in
  the kernel of the full matrix, whose upper component is $Qy$.
\end{proof}

\begin{theorem}[Pseudoinverse completion of the square]
  \label{thm:schur_complement_quadratic_form}
  \lean{SchurComplement.schur_complement_quadratic_form}
  \leanok
  \uses{def:blockMatrix, def:schurComplement}
  If $R$ is positive semi-definite and
  $\ker(R)\subseteq\ker(Q)$, then
  \begin{align}
    \begin{pmatrix}x\\y\end{pmatrix}^{\!\dagger}
      \begin{pmatrix}P&Q\\Q^\dagger&R\end{pmatrix}
      \begin{pmatrix}x\\y\end{pmatrix}
    &=(R^+Q^\dagger x+y)^\dagger R(R^+Q^\dagger x+y)
      +x^\dagger(P-QR^+Q^\dagger)x.
  \end{align}
\end{theorem}

\begin{proof}\leanok
  \uses{thm:support_absorption_of_kernel_inclusion,
    thm:schur_pinv_support_identities}
  The kernel inclusion gives
  $Q\operatorname{supp}(R)=Q$. Expanding the right-hand side and using
  $RR^+=R^+R=\operatorname{supp}(R)$ and $(R^+)^\dagger=R^+$ yields the
  four terms of the block quadratic form.
\end{proof}

\begin{theorem}[Schur-complement positivity from block positivity]
  \label{thm:schur_complement_psd_of_block_psd}
  \lean{SchurComplement.schurComplement_posSemidef_of_blockMatrix_posSemidef}
  \leanok
  \uses{def:blockMatrix, def:schurComplement}
  If $P$ and $R$ are positive semi-definite and
  $\begin{pmatrix}P&Q\\Q^\dagger&R\end{pmatrix}$ is positive
  semi-definite, then $P-QR^+Q^\dagger$ is positive semi-definite.
\end{theorem}

\begin{proof}\leanok
  \uses{thm:block_psd_kernel_inclusion,
    thm:schur_complement_quadratic_form,
    thm:schur_pinv_support_identities}
  Since $P$ and $R^+$ are Hermitian, so are $QR^+Q^\dagger$ and
  $P-QR^+Q^\dagger$.
  By Theorem~\ref{thm:block_psd_kernel_inclusion},
  $\ker(R)\subseteq\ker(Q)$. In
  Theorem~\ref{thm:schur_complement_quadratic_form}, take
  $y=-R^+Q^\dagger x$. The completed-square term vanishes, and positivity of
  the block matrix gives
  $x^\dagger(P-QR^+Q^\dagger)x\ge 0$ for every $x$.
\end{proof}

\begin{theorem}[Block positivity from Schur-complement positivity]
  \label{thm:block_psd_of_schur_complement_psd}
  \lean{SchurComplement.blockMatrix_posSemidef_of_schurComplement_posSemidef}
  \leanok
  \uses{def:blockMatrix, def:schurComplement}
  Suppose that $P$ and $R$ are positive semi-definite,
  $\ker(R)\subseteq\ker(Q)$, and $P-QR^+Q^\dagger$ is positive
  semi-definite. Then
  $\begin{pmatrix}P&Q\\Q^\dagger&R\end{pmatrix}$ is positive
  semi-definite.
\end{theorem}

\begin{proof}\leanok
  \uses{lem:blockMatrix_isHermitian,
    thm:schur_complement_quadratic_form}
  Theorem~\ref{thm:schur_complement_quadratic_form} writes every quadratic
  form of the block matrix as the sum of quadratic forms of $R$ and of
  $P-QR^+Q^\dagger$. Both are nonnegative.
\end{proof}

\begin{theorem}[Block positivity and the pseudoinverse Schur complement]
  \label{thm:block_psd_iff_schur_complement_psd}
  \lean{SchurComplement.blockMatrix_posSemidef_iff}
  \leanok
  \uses{def:blockMatrix, def:schurComplement}
  Let $P$ and $R$ be positive semi-definite. Then
  \begin{align}
    \begin{pmatrix}P&Q\\Q^\dagger&R\end{pmatrix}\ge 0
    \quad\Longleftrightarrow\quad
    \ker(R)\subseteq\ker(Q)
    \ \text{and}\ P-QR^+Q^\dagger\ge 0.
  \end{align}

  \textbf{Scope restriction (Wolf Theorem 5.2, conditions (1) and (2)):}
  This proves exactly the first two conditions of the source theorem. The
  printed contraction condition (3) is false in the singular case, as recorded in
  \texttt{docs/paper-gaps/schur\_complement\_tfae.tex}.
\end{theorem}

\begin{proof}\leanok
  \uses{thm:block_psd_kernel_inclusion,
    thm:schur_complement_psd_of_block_psd,
    thm:block_psd_of_schur_complement_psd}
  The forward implication combines kernel inclusion with positivity of the
  Schur complement. The reverse implication is
  Theorem~\ref{thm:block_psd_of_schur_complement_psd}.
\end{proof}

\begin{theorem}[Pseudoinverse as inverse on the support]
  \label{thm:schur_pinv_eq_support_inverse}
  \lean{SchurComplement.pinv_eq_supportInv}
  \leanok
  \uses{def:pinv, thm:schur_pinv_support_identities}
  If $R$ is positive semi-definite, then its Moore--Penrose pseudoinverse
  agrees with the inverse obtained by inverting $R$ on its support and setting
  it to zero on the kernel.
\end{theorem}

\begin{definition}[Support-normalized off-diagonal block]
  \label{def:schur_contraction}
  \lean{SchurComplement.schurContraction}
  \leanok
  For positive semi-definite $P$ and $R$, define
  \begin{align}
    K=P^{-1/2}QR^{-1/2},
  \end{align}
  where both inverse square roots vanish on the respective kernels.
\end{definition}

\begin{theorem}[Normalized Schur term]
  \label{thm:schur_contraction_mul_adjoint}
  \lean{SchurComplement.schurContraction_mul_conjTranspose}
  \leanok
  \uses{def:schur_contraction, thm:schur_pinv_eq_support_inverse}
  The support-normalized off-diagonal block satisfies
  \begin{align}
    KK^\dagger=P^{-1/2}(QR^+Q^\dagger)P^{-1/2}.
  \end{align}
\end{theorem}

\begin{theorem}[Positivity of the Schur term]
  \label{thm:schur_term_psd}
  \lean{SchurComplement.schurTerm_posSemidef}
  \leanok
  \uses{thm:schur_pinv_eq_support_inverse}
  If $R$ is positive semi-definite, then $QR^+Q^\dagger$ is positive
  semi-definite.
\end{theorem}

\begin{theorem}[Congruence criterion for domination]
  \label{thm:psd_congruence_norm_criterion}
  \lean{Matrix.PosSemidef.sub_smul_posSemidef_iff}
  \leanok
  Let $A$ and $B$ be positive semi-definite, let $c>0$, and suppose that
  $\ker(A)\subseteq\ker(B)$. Then
  \begin{align}
    A-cB\ge0
    \quad\Longleftrightarrow\quad
    c\lVert A^{-1/2}BA^{-1/2}\rVert_\infty\le1,
  \end{align}
  where the inverse square root vanishes on $\ker(A)$.
\end{theorem}

\begin{theorem}[Left support from Schur-complement positivity]
  \label{thm:left_support_of_schur_complement_psd}
  \lean{SchurComplement.ker_conjTranspose_subset_of_schurComplement_posSemidef}
  \leanok
  \uses{def:schurComplement, thm:schur_term_psd,
    thm:support_absorption_of_kernel_inclusion}
  Suppose that $P$ and $R$ are positive semi-definite,
  $\ker(R)\subseteq\ker(Q)$, and $P-QR^+Q^\dagger$ is positive
  semi-definite. Then $\ker(P)\subseteq\ker(Q^\dagger)$.
\end{theorem}

\begin{theorem}[Corrected Schur-complement contraction criterion]
  \label{thm:schur_complement_psd_iff_corrected_contraction}
  \lean{SchurComplement.schurComplement_posSemidef_iff_contraction}
  \leanok
  \uses{def:schurComplement, def:schur_contraction}
  Let $P$ and $R$ be positive semi-definite and suppose that
  $\ker(R)\subseteq\ker(Q)$. Then
  \begin{align}
    P-QR^+Q^\dagger\ge0
    \quad\Longleftrightarrow\quad
    \ker(P)\subseteq\ker(Q^\dagger)
    \ \text{and}\ \lVert P^{-1/2}QR^{-1/2}\rVert_\infty\le1.
  \end{align}

  \textbf{Local fix (Wolf Theorem 5.2, condition (3)):}
  The left-support condition is necessary and is absent from the printed
  statement. See
  \texttt{docs/paper-gaps/schur\_complement\_tfae.tex}.
\end{theorem}

\begin{proof}\leanok
  \uses{thm:psd_congruence_norm_criterion,
    thm:schur_contraction_mul_adjoint, thm:schur_term_psd,
    thm:left_support_of_schur_complement_psd}
  Put $B=QR^+Q^\dagger$ and $K=P^{-1/2}QR^{-1/2}$. The two support conditions
  allow congruence by $P^{-1/2}$ to identify $P-B\ge0$ with
  $\lVert P^{-1/2}BP^{-1/2}\rVert_\infty\le1$. By
  Theorem~\ref{thm:schur_contraction_mul_adjoint}, the matrix inside this norm
  is $KK^\dagger$, whose norm is $\lVert K\rVert_\infty^2$.
\end{proof}

\begin{theorem}[Counterexample to Wolf's printed condition (3)]
  \label{thm:wolf_schur_condition_three_counterexample}
  \lean{SchurComplement.wolf_condition_three_not_sufficient}
  \leanok
  \uses{def:blockMatrix, def:schur_contraction,
    thm:block_quadratic_form}
  For the scalar blocks $P=0$ and $Q=R=1$, the printed kernel condition and
  contraction bound hold, but
  $\begin{pmatrix}P&Q\\Q^\dagger&R\end{pmatrix}$ is not positive
  semi-definite.
\end{theorem}

\begin{proof}\leanok
  The normalized off-diagonal block vanishes because $P^{-1/2}=0$. However,
  the quadratic form of $\begin{pmatrix}0&1\\1&1\end{pmatrix}$ at $(1,-1)$
  equals $-1$.
\end{proof}

\begin{theorem}[Corrected block contraction criterion]
  \label{thm:block_psd_iff_corrected_contraction}
  \lean{SchurComplement.blockMatrix_posSemidef_iff_contraction}
  \leanok
  \uses{def:blockMatrix, def:schur_contraction}
  Let $P$ and $R$ be positive semi-definite. Then
  \begin{align}
    \begin{pmatrix}P&Q\\Q^\dagger&R\end{pmatrix}\ge0
    \quad\Longleftrightarrow\quad
    &\ker(R)\subseteq\ker(Q),\\
    &\ker(P)\subseteq\ker(Q^\dagger),\\
    &\lVert P^{-1/2}QR^{-1/2}\rVert_\infty\le1.
  \end{align}

  \textbf{Local fix (Wolf Theorem 5.2, condition (3)):}
  This is the source statement with its missing left-support condition restored.
  See \texttt{docs/paper-gaps/schur\_complement\_tfae.tex}.
\end{theorem}

\begin{proof}\leanok
  \uses{thm:block_psd_kernel_inclusion,
    thm:schur_complement_psd_of_block_psd,
    thm:block_psd_of_schur_complement_psd,
    thm:schur_complement_psd_iff_corrected_contraction}
  Combine the pseudoinverse Schur-complement criterion with the corrected
  contraction criterion.
\end{proof}
